\begin{lemma}[Archimedean Integer Part Lemma]
\label{lem:archimedean-integer-part-lemma}
For every $x\in\mathbb{R}$, there exists a unique integer $m\in\mathbb{Z}$
such that
\[
  m\leq x<m+1.
\]
Equivalently,
\[
  (\forall x\in\mathbb{R})(\exists!m\in\mathbb{Z})(m\leq x<m+1).
\]

\smallskip
\noindent
\hyperref[prf:archimedean-integer-part-lemma]{\textit{Go to proof.}}
\end{lemma}
\end{lemmabox}

\begin{remark*}[Standard quantified statement]
\[
\forall x\in\mathbb{R}\;\exists m\in\mathbb{Z}\;(m-1\le x<m).
\]
\end{remark*}

\begin{remark*}[Predicate reading]
\[
m=\lfloor x\rfloor+1
\Longleftrightarrow
m\in\mathbb{Z}\land m-1\le x<m.
\]
\end{remark*}

\begin{remark*}[Negated quantified statement]
\[
\exists x\in\mathbb{R}\;\forall m\in\mathbb{Z}\;(x<m-1\lor m\le x).
\]
\end{remark*}

\begin{remark*}[Negation predicate reading]
\[
\neg\exists m\in\mathbb{Z}\;(m=\lfloor x\rfloor+1)
\Longleftrightarrow
\forall m\in\mathbb{Z}\;(x<m-1\lor m\le x).
\]
\end{remark*}

\begin{remark*}[Failure modes]
\begin{description}
  \item[Exposition.]
The lemma fails only if some real number lies in no half-open unit interval
$[m-1,m)$ with integer right endpoint.


  \item[\exists-condition.]
  \[
\exists x\in\mathbb{R}\;\forall m\in\mathbb{Z}\;(x<m-1\lor m\le x).
\]
\[
\neg\exists m\in\mathbb{Z}\;(m=\lfloor x\rfloor+1)
\Longleftrightarrow
\forall m\in\mathbb{Z}\;(x<m-1\lor m\le x).
\]
\end{description}
\end{remark*}

\begin{remark*}[Failure modes]
\begin{description}
  \item[Exposition.]
\[
\underbrace{x<m-1}_{\text{right endpoint too high}}
\qquad\text{or}\qquad
\underbrace{m\le x}_{\text{right endpoint not above }x}.
\]


  \item[\exists-condition.]
  \[
\exists x\in\mathbb{R}\;\forall m\in\mathbb{Z}\;(x<m-1\lor m\le x).
\]
\[
\neg\exists m\in\mathbb{Z}\;(m=\lfloor x\rfloor+1)
\Longleftrightarrow
\forall m\in\mathbb{Z}\;(x<m-1\lor m\le x).
\]
\end{description}
\end{remark*}

\begin{remark*}[Interpretation]
This is the integer-part lemma with the endpoint shifted by one.  It produces
the first integer strictly above $x$ when the unit interval is written with its
integer endpoint on the right.
\end{remark*}

\begin{dependencies}
\begin{itemize}
\item \hyperref[def:reals]{Real Numbers}.
\item \hyperref[def:integers]{Integers}.
\item \hyperref[lem:integer-part-lemma]{Integer Part Lemma}.
\end{itemize}
\end{dependencies}



